\documentclass[aspectratio=169]{beamer}

\usepackage[utf8]{inputenc}
\usepackage[T1]{fontenc}
\usepackage{lmodern}
\usepackage{helvet}
\renewcommand{\familydefault}{\sfdefault}
\usepackage{graphicx}
\usepackage{booktabs}
\usepackage{array}
\usepackage{xspace}
\usepackage{amsmath}
\usepackage{tabularx}

\graphicspath{{../outputs/figures/}}

% ---------------------------------------------------------------------------
% Cores (mesma paleta categórica usada nas figuras geradas por
% scripts/08_make_presentation_assets.py)
% ---------------------------------------------------------------------------
\definecolor{accentblue}{HTML}{2A78D6}
\definecolor{accentorange}{HTML}{EB6834}
\definecolor{accentred}{HTML}{E34948}
\definecolor{inkprimary}{HTML}{0B0B0B}
\definecolor{inksecondary}{HTML}{52514E}
\definecolor{gridgray}{HTML}{E1E0D9}

\setbeamercolor{normal text}{fg=inkprimary,bg=white}
\setbeamercolor{structure}{fg=accentblue}
\setbeamercolor{title}{fg=accentblue}
\setbeamercolor{frametitle}{fg=inkprimary,bg=white}
\setbeamercolor{footline}{fg=inksecondary,bg=white}
\setbeamercolor{block title}{fg=white,bg=accentblue}
\setbeamercolor{block body}{fg=inkprimary,bg=white}
\setbeamercolor{alerted text}{fg=accentorange}
\setbeamercolor{itemize item}{fg=accentblue}
\setbeamercolor{itemize subitem}{fg=accentorange}

\setbeamertemplate{navigation symbols}{}
\setbeamertemplate{itemize items}[circle]
\setbeamertemplate{frametitle}{%
  \vspace{0.25cm}
  {\usebeamercolor[fg]{frametitle}\Large\bfseries\insertframetitle}\par
  \vspace{2pt}
  {\color{accentblue}\rule{\linewidth}{1.4pt}}
  \vspace{4pt}
}
\setbeamertemplate{footline}{%
  \leavevmode
  \hbox{\begin{beamercolorbox}[wd=\paperwidth,ht=2.4ex,dp=1.2ex,leftskip=1em,rightskip=1em]{footline}
  {\scriptsize Formalidade, g\^enero, ra\c{c}a/cor e renda em Roraima}
  \hfill
  {\scriptsize\insertframenumber\,/\,\inserttotalframenumber}
  \end{beamercolorbox}}
}

\newcommand{\highlight}[1]{\textcolor{accentblue}{\textbf{#1}}}
\newcommand{\negativo}[1]{\textcolor{accentred}{\textbf{#1}}}

\input{gerado/numeros.tex}

\title[Formalidade, g\^enero, ra\c{c}a/cor e renda em RR]{Formalidade, g\^enero, ra\c{c}a/cor e renda\\no mercado de trabalho de Roraima}
\subtitle{Diferenciais condicionais de rendimento com microdados da PNAD Cont\'inua}
\author{Yuri Cesar Silva}
\institute{Resultados preliminares de pesquisa}
\date{Agosto de 2026}

\begin{document}

\begin{frame}[plain]
  \titlepage
\end{frame}

\begin{frame}{Sum\'ario}
  \tableofcontents
\end{frame}

% =============================================================================
\section{Introdu\c{c}\~ao}
% =============================================================================

\begin{frame}{O problema}
  \begin{itemize}
    \item Roraima tem um mercado de trabalho pequeno, com forte presen\c{c}a de informalidade e composi\c{c}\~ao demogr\'afica marcada por fluxo migrat\'orio (fronteira com a Venezuela).
    \item Diferenciais de rendimento por formalidade, g\^enero e ra\c{c}a/cor s\~ao amplamente documentados para o Brasil como um todo, mas pouco explorados especificamente para Roraima.
    \item \highlight{Pergunta central:} mantidas constantes ocupa\c{c}\~ao, atividade econ\^omica, escolaridade, idade, jornada e per\'iodo, persistem diferenciais de rendimento associados \`a formalidade, g\^enero e ra\c{c}a/cor no mercado de trabalho de Roraima?
  \end{itemize}
\end{frame}

\begin{frame}{Por que isso importa}
  \begin{itemize}
    \item Formalidade: acesso a direitos trabalhistas, previd\^encia e prote\c{c}\~ao social \'e mediado pelo v\'inculo formal.
    \item G\^enero e ra\c{c}a/cor: diferenciais persistentes mesmo ap\'os controle por escolaridade e ocupa\c{c}\~ao s\~ao evid\^encia consistente com discrimina\c{c}\~ao ou com barreiras de acesso a determinados postos de trabalho.
    \item Roraima \'e um caso pouco estudado: economia pequena, fortemente dependente do setor p\'ublico, com choque migrat\'orio recente relevante para o mercado de trabalho local.
    \item Entender a composi\c{c}\~ao desses diferenciais informa pol\'iticas de formaliza\c{c}\~ao, equidade salarial e igualdade racial no estado.
  \end{itemize}
\end{frame}

\begin{frame}{Um primeiro olhar: gaps brutos (sem nenhum controle)}
  \begin{center}
    \includegraphics[width=0.95\linewidth,height=0.78\textheight,keepaspectratio]{motivacao_gaps_brutos.pdf}
  \end{center}
  \vspace{2mm}
  {\footnotesize Barra azul = grupo de refer\^encia; barra laranja = grupo de compara\c{c}\~ao. M\'edias ponderadas pelo peso amostral, sem nenhum controle.}
\end{frame}

\begin{frame}{O quebra-cabe\c{c}a de g\^enero}
  \begin{itemize}
    \item O gap \highlight{bruto} de renda por hora entre homens e mulheres em Roraima \'e pequeno: \BrutoGapMulherPct (\BrutoMulherGrupoRef\ vs.\ \BrutoMulherGrupoComp).
    \item Isso poderia sugerir, \`a primeira vista, aus\^encia de desigualdade de g\^enero no mercado de trabalho local.
    \item \alert{Mas} essa compara\c{c}\~ao ignora que homens e mulheres est\~ao distribu\'idos de forma desigual entre ocupa\c{c}\~oes, setores, escolaridade e jornada.
    \item Ao longo desta apresenta\c{c}\~ao, mostramos que o gap \highlight{condicional} (mantendo ocupa\c{c}\~ao, setor, escolaridade, idade e per\'iodo fixos) \'e substancialmente maior: \negativo{\PctMulherHora} na renda por hora.
    \item Esse \'e um exemplo cl\'assico de por que compara\c{c}\~oes brutas podem esconder desigualdades reais.
  \end{itemize}
\end{frame}

\begin{frame}{Objetivos}
  \begin{enumerate}
    \item Estimar o pr\^emio salarial da formalidade em Roraima.
    \item Estimar os diferenciais condicionais de g\^enero e ra\c{c}a/cor, controlando por composi\c{c}\~ao ocupacional e setorial.
    \item Testar se a formaliza\c{c}\~ao atenua ou amplia os diferenciais de g\^enero e ra\c{c}a/cor (termos de intera\c{c}\~ao).
    \item Decompor o diferencial de renda mensal em efeito-pre\c{c}o (renda por hora) e efeito-jornada (horas trabalhadas).
    \item Reportar os resultados em termos economicamente interpret\'aveis: \% e R\$/horas.
  \end{enumerate}
\end{frame}

% =============================================================================
\section{Revis\~ao de literatura}
% =============================================================================

\begin{frame}{Base te\'orica: equa\c{c}\~ao de rendimentos}
  \begin{itemize}
    \item A especifica\c{c}\~ao segue a tradi\c{c}\~ao da \highlight{equa\c{c}\~ao de rendimentos mincer\-iana}: log do rendimento como fun\c{c}\~ao de capital humano (escolaridade, experi\^encia/idade) e caracter\'isticas do posto de trabalho.
    \begin{itemize}
      \item Mincer, J. (1974). \textit{Schooling, Experience, and Earnings}. NBER.
      \item Card, D. (1999). ``The Causal Effect of Education on Earnings''. \textit{Handbook of Labor Economics}, v.\ 3A.
    \end{itemize}
    \item A extens\~ao natural inclui indicadores de \highlight{tratamento} (formalidade, g\^enero, ra\c{c}a/cor) ao lado dos controles de capital humano -- pr\'atica padr\~ao na literatura de diferenciais salariais (Oaxaca-Blinder e extens\~oes com efeitos fixos).
  \end{itemize}
\end{frame}

\begin{frame}{Informalidade no mercado de trabalho}
  \begin{itemize}
    \item Ulyssea, G. (2018). ``Firms, Informality, and Development: Theory and Evidence from Brazil''. \textit{American Economic Review}, 108(8).
    \begin{itemize}
      \item Modela a informalidade como resposta de firmas a custos de regula\c{c}\~ao e fiscaliza\c{c}\~ao, com implica\c{c}\~oes para produtividade e aloca\c{c}\~ao de trabalhadores.
    \end{itemize}
    \item Achado emp\'irico recorrente na literatura brasileira: trabalhadores formais recebem pr\^emio salarial mesmo ap\'os controle por escolaridade e ocupa\c{c}\~ao, consistente com segmenta\c{c}\~ao (n\~ao apenas sele\c{c}\~ao) do mercado de trabalho.
  \end{itemize}
\end{frame}

\begin{frame}{G\^enero e ra\c{c}a/cor no mercado de trabalho}
  \begin{itemize}
    \item Blau, F. \& Kahn, L. (2017). ``The Gender Wage Gap: Extent, Trends, and Explanations''. \textit{Journal of Economic Literature}, 55(3).
    \begin{itemize}
      \item Parte do gap de g\^enero \'e explicada por diferen\c{c}as observ\'aveis (ocupa\c{c}\~ao, jornada); uma parcela relevante permanece n\~ao explicada mesmo em especifica\c{c}\~oes ricas em controles.
    \end{itemize}
    \item Literatura brasileira sobre desigualdade racial de rendimentos documenta diferenciais persistentes entre brancos e pretos/pardos mesmo ap\'os controle por escolaridade, regi\~ao e setor formal/informal.
    \item \alert{Nota:} esta \'e uma revis\~ao inicial. Citações espec\'ificas para o contexto de Roraima/regi\~ao Norte, fronteira e fluxo migrat\'orio venezuelano ainda precisam ser levantadas e incorporadas na pr\'oxima rodada.
  \end{itemize}
\end{frame}

% =============================================================================
\section{Dados e metodologia}
% =============================================================================

\begin{frame}{Fonte de dados e amostra}
  \begin{itemize}
    \item \highlight{Fonte:} microdados trimestrais oficiais da PNAD Cont\'inua (IBGE).
    \item \highlight{Recorte:} indiv\'iduos ocupados em Roraima (UF~14), 4\textordmasculine{} trimestre de cada ano, 2016T4--2025T4 -- ``fotografias'' de fim de per\'iodo (empilhamento de todos os trimestres fica como robustez futura).
    \item \highlight{Funil amostral (n\~ao ponderado):}
    \begin{itemize}
      \item Total de entrevistados: \NTotal
      \item Popula\c{c}\~ao de 14 anos ou mais: \NQuatorzeMais
      \item Ocupados: \NOcupados
      \item Ocupados com rendimento v\'alido: \NRendaValida
    \end{itemize}
    \item Estima\c{c}\~ao: WLS com peso amostral (\texttt{V1028}) e erros-padr\~ao clusterizados por UPA.
  \end{itemize}
\end{frame}

\begin{frame}{Vari\'aveis}
  \begin{columns}[T]
    \begin{column}{0.5\linewidth}
      \textbf{Dependentes}
      \begin{itemize}
        \item $\ln(\text{renda mensal real})$
        \item $\ln(\text{renda por hora real})$
      \end{itemize}
      \vspace{2mm}
      \textbf{Tratamentos}
      \begin{itemize}
        \item \texttt{formal}: posi\c{c}\~ao na ocupa\c{c}\~ao/carteira assinada, com CNPJ como crit\'erio auxiliar para empregadores/conta-pr\'opria
        \item \texttt{mulher}: sexo feminino
        \item \texttt{preto\_pardo}: pretos ou pardos (ref.: brancos, amarelos ou ind\'igenas)
      \end{itemize}
    \end{column}
    \begin{column}{0.5\linewidth}
      \textbf{Controles} (agrupados como ``controles'' nas tabelas)
      \begin{itemize}
        \item Idade e idade\textsuperscript{2}
        \item Escolaridade (7 categorias)
        \item Grande grupo ocupacional (11 categorias)
        \item Grande grupo de atividade (12 categorias)
        \item Efeitos fixos de ano-trimestre (10 per\'iodos)
      \end{itemize}
    \end{column}
  \end{columns}
\end{frame}

\begin{frame}{Renda real: deflacionamento}
  \begin{itemize}
    \item Rendimento nominal (\texttt{VD4016}) deflacionado pelo \highlight{deflator oficial trimestral da PNAD Cont\'inua} (IBGE), espec\'ifico por UF e trimestre.
    \item Renda m\'edia real na amostra: \MediaMensal/m\^es e \MediaHora/hora; jornada m\'edia \MediaHoras\ horas/semana.
    \item \alert{Nota metodol\'ogica:} como o modelo j\'a satura efeitos fixos de ano-trimestre, os coeficientes de \texttt{formal}, \texttt{mulher} e \texttt{preto\_pardo} s\~ao numericamente id\^enticos entre renda nominal e real -- o deflator (constante dentro de cada per\'iodo) \'e absorvido pelas pr\'oprias dummies de per\'iodo. A defla\c{c}\~ao continua sendo necess\'aria para descrever a evolu\c{c}\~ao do \emph{n\'ivel} de renda ao longo do tempo, o que fazemos nas convers\~oes para R\$ adiante.
  \end{itemize}
\end{frame}

\begin{frame}{Especifica\c{c}\~ao do modelo-base}
  \small
  \begin{block}{Modelo principal (WLS, erros clusterizados por UPA)}
    \setlength{\abovedisplayskip}{2pt}
    \setlength{\belowdisplayskip}{2pt}
    \begin{align*}
      \ln(w_{it}) = {} & \beta_1\, formal_{it} + \beta_2\, mulher_i + \beta_3\, preto\_pardo_i \\
                        & + \beta_4\, (formal_{it}\times mulher_i) + \beta_5\, (formal_{it}\times preto\_pardo_i) \\
                        & + X_{it}\,\gamma + FE_{ocupacao} + FE_{atividade} + FE_{periodo} + \varepsilon_{it}
    \end{align*}
  \end{block}
  \begin{itemize}
    \item $w_{it}$: rendimento mensal real ou rendimento por hora real (dois modelos separados).
    \item Estimamos tamb\'em uma vers\~ao \highlight{sem controles} (apenas os tratamentos e suas intera\c{c}\~oes) para mostrar o quanto os coeficientes mudam ao controlar por composi\c{c}\~ao -- crit\'erio de robustez principal desta apresenta\c{c}\~ao.
    \item Leitura causal exige cautela: a sele\c{c}\~ao para o emprego formal \'e potencialmente end\'ogena.
  \end{itemize}
\end{frame}

% =============================================================================
\section{Resultados}
% =============================================================================

\begin{frame}{Resultado 1 -- Renda mensal real}
  \begin{center}
    \scriptsize
    \renewcommand{\arraystretch}{0.82}
    \setlength{\tabcolsep}{5pt}
    \input{gerado/tab_mensal.tex}
  \end{center}
  {\scriptsize Coeficientes de $\ln(\text{renda mensal real})$. Erros-padr\~ao clusterizados por UPA entre par\^enteses. $^{*}p<0{,}10$; $^{**}p<0{,}05$; $^{***}p<0{,}01$. Apenas as vari\'aveis de interesse s\~ao exibidas; a coluna (2) inclui todos os controles listados na metodologia.}
\end{frame}

\begin{frame}{Resultado 2 -- Renda por hora real}
  \begin{center}
    \scriptsize
    \renewcommand{\arraystretch}{0.82}
    \setlength{\tabcolsep}{5pt}
    \input{gerado/tab_hora.tex}
  \end{center}
  {\scriptsize Coeficientes de $\ln(\text{renda por hora real})$. Erros-padr\~ao clusterizados por UPA entre par\^enteses. $^{*}p<0{,}10$; $^{**}p<0{,}05$; $^{***}p<0{,}01$.}
\end{frame}

\begin{frame}{Robustez: sem controles vs.\ com controles}
  \begin{center}
    \includegraphics[width=0.95\linewidth,height=0.74\textheight,keepaspectratio]{coef_forest.pdf}
  \end{center}
\end{frame}

\begin{frame}{O que a Figura anterior mostra}
  \begin{itemize}
    \item \highlight{Formal:} o pr\^emio bruto (\BrutoGapFormalPct\ na renda por hora, comparando m\'edias) \'e superestimado sem controles -- cai de forma acentuada ao controlar por ocupa\c{c}\~ao/atividade/escolaridade, mas continua grande e significativo (\PctFormalHora\ na renda por hora).
    \item \highlight{Mulher:} o sinal at\'e \highlight{inverte} entre as especifica\c{c}\~oes na renda por hora -- sem controles o coeficiente \'e pr\'oximo de zero/positivo; com controles, torna-se \negativo{\PctMulherHora}. Sinal de que mulheres est\~ao concentradas em ocupa\c{c}\~oes/setores de rendimento mais alto, o que mascara o gap salarial dentro da mesma ocupa\c{c}\~ao.
    \item \highlight{Preto/pardo:} o gap \'e negativo e significativo nas duas especifica\c{c}\~oes, com magnitude menor ao controlar por composi\c{c}\~ao.
  \end{itemize}
\end{frame}

% =============================================================================
\section{Mecanismo: pre\c{c}o do trabalho vs.\ jornada}
% =============================================================================

\begin{frame}{Decompondo o diferencial de renda mensal}
  \begin{itemize}
    \item O diferencial de renda \highlight{mensal} mistura dois canais: quanto se ganha por hora (pre\c{c}o) e quantas horas se trabalha (quantidade).
    \item Estrat\'egia: comparar \emph{(i)} renda mensal sem controlar horas, \emph{(ii)} renda mensal controlando por horas, \emph{(iii)} renda por hora, e \emph{(iv)} horas semanais como vari\'avel dependente -- todos na mesma amostra (renda e horas v\'alidas, $N=$\NJornada).
  \end{itemize}
  \begin{center}
    \footnotesize
    \input{gerado/tab_mecanismo.tex}
  \end{center}
\end{frame}

\begin{frame}{Mecanismo de jornada -- visual}
  \begin{center}
    \includegraphics[width=0.95\linewidth,height=0.78\textheight,keepaspectratio]{mecanismo_jornada.pdf}
  \end{center}
\end{frame}

\begin{frame}{Leitura do mecanismo}
  \begin{itemize}
    \item \highlight{Mulheres} trabalham, em m\'edia, \HorasMulher\ horas/semana em rela\c{c}\~ao a homens compar\'aveis -- efeito-jornada relevante. O gap mensal cai de \PctMulherMensalSemHorasJornada\ para \PctMulherMensalComHorasJornada\ ao controlar horas, e o gap por hora \'e \PctMulherHoraJornada. Ou seja: o diferencial de g\^enero \'e \highlight{misto} -- parte pre\c{c}o, parte jornada mais curta.
    \item \highlight{Formal:} trabalhadores formais trabalham \emph{mais} horas (\HorasFormal\ h/sem) e ainda assim ganham mais por hora -- o pr\^emio de formalidade \highlight{n\~ao} \'e explicado por jornada mais curta.
    \item \highlight{Preto/pardo:} efeito-jornada pequeno (\HorasPretoPardo\ h/sem); o gap \'e majoritariamente efeito-pre\c{c}o.
  \end{itemize}
\end{frame}

% =============================================================================
\section{Discuss\~ao: quanto isso vale em R\$?}
% =============================================================================

\begin{frame}{Convertendo \% em R\$ (na renda m\'edia real da amostra)}
  \begin{center}
    \includegraphics[width=0.9\linewidth,height=0.62\textheight,keepaspectratio]{gap_reais.pdf}
  \end{center}
  {\footnotesize Convers\~ao ilustrativa: efeito percentual do modelo com controles aplicado \`a renda m\'edia real da amostra (\MediaMensal/m\^es; \MediaHora/hora). N\~ao \'e uma previs\~ao para um indiv\'iduo espec\'ifico, mas uma refer\^encia de ordem de grandeza.}
\end{frame}

\begin{frame}{Discuss\~ao -- Formalidade}
  \begin{itemize}
    \item Trabalhadores formais ganham \highlight{\PctFormalMensal} a mais por m\^es (\highlight{\ReaisFormalMensal}) e \highlight{\PctFormalHora} a mais por hora (\highlight{\ReaisFormalHora}) do que informais compar\'aveis.
    \item \'E o maior diferencial estimado no estudo -- maior que g\^enero ou ra\c{c}a/cor.
    \item A intera\c{c}\~ao formal$\times$mulher \'e positiva e significativa na renda mensal (\PctFormalMulherMensal, $p=\PValorFormalMulherMensal$): a formaliza\c{c}\~ao \highlight{atenua} o gap mensal de g\^enero, plausivelmente por padronizar jornada. Na renda por hora a intera\c{c}\~ao n\~ao \'e significativa ($p=\PValorFormalMulherHora$).
    \item A intera\c{c}\~ao formal$\times$preto/pardo n\~ao \'e significativa em nenhum dos dois modelos ($p=\PValorFormalPretoPardoMensal$ mensal; $p=\PValorFormalPretoPardoHora$ hora).
  \end{itemize}
\end{frame}

\begin{frame}{Discuss\~ao -- G\^enero}
  \begin{itemize}
    \item Mulheres ganham \negativo{\PctMulherMensal} a menos por m\^es (\negativo{\ReaisMulherMensal}) e \negativo{\PctMulherHora} a menos por hora (\negativo{\ReaisMulherHora}) do que homens compar\'aveis, mantendo ocupa\c{c}\~ao, atividade, escolaridade, idade e per\'iodo fixos.
    \item O gap \highlight{bruto} (sem controles) era pequeno (\BrutoGapMulherPct) -- a desigualdade s\'o aparece claramente quando comparamos mulheres e homens \emph{dentro da mesma ocupa\c{c}\~ao e setor}.
    \item Parte do gap mensal \'e jornada mais curta (\HorasMulher\ h/sem); parte \'e pre\c{c}o menor por hora trabalhada.
  \end{itemize}
\end{frame}

\begin{frame}{Discuss\~ao -- Ra\c{c}a/cor}
  \begin{itemize}
    \item Pessoas pretas ou pardas ganham \negativo{\PctPretoPardoMensal} a menos por m\^es (\negativo{\ReaisPretoPardoMensal}) e \negativo{\PctPretoPardoHora} a menos por hora (\negativo{\ReaisPretoPardoHora}) do que o grupo de refer\^encia, mantendo os mesmos controles fixos.
    \item Diferente do g\^enero, o gap racial \'e vis\'ivel mesmo na compara\c{c}\~ao bruta (\BrutoGapPretoPardoPct) e permanece ap\'os controles -- embora de magnitude menor que os gaps de formalidade e g\^enero.
    \item Efeito-jornada pequeno: o gap racial \'e majoritariamente efeito-pre\c{c}o (renda por hora), n\~ao diferen\c{c}a de horas trabalhadas.
  \end{itemize}
\end{frame}

\begin{frame}{Cautelas de interpreta\c{c}\~ao}
  \begin{itemize}
    \item Estes s\~ao \highlight{diferenciais condicionais}, n\~ao efeitos causais: controlamos por ocupa\c{c}\~ao, atividade, escolaridade, idade e per\'iodo, mas a sele\c{c}\~ao para emprego formal, para ocupa\c{c}\~ao e para jornada continua potencialmente end\'ogena.
    \item A convers\~ao para R\$ usa a renda m\'edia da amostra como refer\^encia -- \'e ilustrativa, n\~ao \'e uma previs\~ao individual.
    \item Amostra restrita ao 4\textordmasculine\ trimestre de cada ano (fotografias de fim de per\'iodo); robustez com todos os trimestres empilhados ainda n\~ao foi feita.
  \end{itemize}
\end{frame}

% =============================================================================
\section{Conclus\~ao}
% =============================================================================

\begin{frame}{S\'intese dos achados}
  \begin{enumerate}
    \item Pr\^emio de formalidade \'e grande e robusto (\PctFormalHora\ na renda/hora) -- n\~ao explicado por jornada.
    \item Gap de g\^enero \'e pequeno na compara\c{c}\~ao bruta mas grande e significativo quando controlamos composi\c{c}\~ao (\PctMulherHora\ na renda/hora); \'e parcialmente explicado por jornada mais curta.
    \item Gap racial \'e menor em magnitude, vis\'ivel mesmo sem controles, e majoritariamente efeito-pre\c{c}o.
    \item Formaliza\c{c}\~ao atenua o gap \emph{mensal} de g\^enero, mas n\~ao o gap racial.
  \end{enumerate}
\end{frame}

\begin{frame}{Pr\'oximos passos}
  \begin{itemize}
    \item Aprofundar a revis\~ao de literatura, com foco em: mercado de trabalho na regi\~ao Norte/fronteira, impacto do fluxo migrat\'orio venezuelano, e estudos raciais/g\^enero espec\'ificos para o Brasil.
    \item R\'eplica das estimativas principais em R com o pacote \texttt{survey}, incorporando desenho amostral completo (estratos e conglomerados).
    \item Ap\^endice de robustez: todos os trimestres empilhados (n\~ao s\'o o 4\textordmasculine); agrega\c{c}\~oes ocupacionais/setoriais alternativas.
    \item Investigar mecanismos adicionais do gap de g\^enero (tipo de ocupa\c{c}\~ao, trabalho dom\'estico/cuidado n\~ao remunerado como hip\'otese para a diferen\c{c}a de horas).
  \end{itemize}
\end{frame}

\begin{frame}{Refer\^encias}
  \footnotesize
  \begin{itemize}
    \item Blau, F.\ D. \& Kahn, L.\ M.\ (2017). The Gender Wage Gap: Extent, Trends, and Explanations. \textit{Journal of Economic Literature}, 55(3), 789--865.
    \item Card, D.\ (1999). The Causal Effect of Education on Earnings. In \textit{Handbook of Labor Economics}, v.\ 3A. Elsevier.
    \item IBGE (2026). Pesquisa Nacional por Amostra de Domic\'ilios Cont\'inua -- Microdados trimestrais e deflatores oficiais.
    \item Mincer, J.\ (1974). \textit{Schooling, Experience, and Earnings}. National Bureau of Economic Research.
    \item Ulyssea, G.\ (2018). Firms, Informality, and Development: Theory and Evidence from Brazil. \textit{American Economic Review}, 108(8), 2015--2047.
  \end{itemize}
  \vspace{2mm}
  {\scriptsize Revis\~ao de literatura preliminar -- refer\^encias adicionais sobre g\^enero/ra\c{c}a no Brasil e mercado de trabalho na regi\~ao Norte ser\~ao incorporadas na pr\'oxima vers\~ao.}
\end{frame}

\begin{frame}[plain]
  \centering
  {\Huge\color{accentblue}\bfseries Obrigado!}
  \vspace{6mm}

  {\large Coment\'arios e sugest\~oes s\~ao bem-vindos.}
\end{frame}

\end{document}
